% Autorska priprava k Realisticky.cz 020203; ulohy jsou autorske a obmenene.
% Zamerne nenahrazuje vyklad presnymi kopiemi cizich textu/prikladu.
\documentclass[11pt,a4paper]{article}
\input{../../sablony/pripravy/priprava.sty}
\usepackage{tikz}
\nastavPripravu{1. ročník}{Rovnice a nerovnice}{74}{Lineární rovnice I: závorky a diskuse}
\begin{document}
\blok{Návaznost a motivace}
Přejdeme od pravidel ke skutečnému řešení. Rovnání obou stran má zůstat logické: odstranit závorky, seskupit neznámé a konstanty, zkontrolovat výsledek.
\begin{definiceBox}
\textbf{Lineární rovnice o jedné neznámé}. Rovnice upravitelná ekvivalentními kroky na tvar $ax+b=0$, kde $a,b\in\mathbb R$. Při $a\ne0$ má právě jeden kořen; při $a=0$ rozhoduje člen $b$.
\end{definiceBox}
\blok{Princip a odvození}
Postup pro $ax+b=0$: je-li $a\ne0$, pak $x=-b/a$; je-li $a=0$ a $b\ne0$, řešení neexistuje; je-li $a=b=0$, je řešením každý přípustný argument.
\blok{Řešené příklady: od jednoduchého ke složitějšímu}
\begin{prikladReseniBox}[title={Příklad 1}]
Vyřešte $3x+7=22$.
\tcblower
\begin{align*}
3x+7&=22 \\
 3x&=15 \\
 x&=5
\end{align*}

Zkouška: $3\cdot5+7=22$.
\end{prikladReseniBox}
\begin{prikladReseniBox}[title={Příklad 2}]
Vyřešte $4-2(x-1)=3x+1$.
\tcblower
\begin{align*}
4-2x+2&=3x+1 \\
 6-2x&=3x+1 \\
 5&=5x \\
 x&=1
\end{align*}
\end{prikladReseniBox}
\poznamka{Často se objeví záměna $-2(x-1)$ za $-2x-2$. Opravuj znaménka v závorkách před sbíráním členů.}
\blok{Samostatné procvičování}
\begin{zadaniReseniBox}[title={Úloha 1}]
Vyřešte $2x-5=-9$.
\tcblower
\begin{align*}
2x&=-4 \\
 x&=-2 \\
 K&=\{-2\}
\end{align*}
\end{zadaniReseniBox}
\begin{zadaniReseniBox}[title={Úloha 2}]
Vyřešte $3(x+2)=2x+9$.
\tcblower
\begin{align*}
3x+6&=2x+9 \\
 x&=3 \\
 K&=\{3\}
\end{align*}
\end{zadaniReseniBox}
\begin{zadaniReseniBox}[title={Úloha 3}]
Vyřešte $5x+1=5x-2$.
\tcblower
\begin{align*}
5x+1&=5x-2 \\
 1&=-2 \\
 K&=\varnothing
\end{align*}
\end{zadaniReseniBox}
\blok{Volitelná rezerva}
\begin{zadaniReseniBox}[title={Úloha 1}]
Pro jaké $x$ platí $2(x-1)+3=2x+1$?
\tcblower
\begin{align*}
2x-2+3&=2x+1 \\
 2x+1&=2x+1 \\
 K&=\mathbb R
\end{align*}
\end{zadaniReseniBox}
\blok{Domácí úkol – pět rovnic od jednodušší po těžší}
\begin{zadaniReseniBox}
\begin{enumerate}[label=\alph*)]
\item $7x-4=3x+8$.
\item $3(x-2)=2x+5$.
\item $5-2(3x-4)=3(x-2)$.
\item $4\{2x-[3-(x-2)]\}=12$.
\item $6\{2x-3[4-(x-1)]\}+5=5x-10$.
\end{enumerate}
\tcblower
\begin{enumerate}[label=\alph*)]
\item \(\begin{aligned}7x-4&=3x+8\\4x&=12\\x&=3.\end{aligned}\)
\item \(\begin{aligned}3x-6&=2x+5\\3x-2x&=11\\x&=11.\end{aligned}\)
\item \(\begin{aligned}13-6x&=3x-6\\19&=9x\\x&=\tfrac{19}{9}.\end{aligned}\)
\item \(\begin{aligned}4\{2x-(5-x)\}&=12\\4(3x-5)&=12\\12x-20&=12\\x&=\tfrac83.\end{aligned}\)
\item \(\begin{aligned}6\{2x-3(5-x)\}+5&=5x-10\\6(5x-15)+5&=5x-10\\30x-85&=5x-10\\25x&=75\\x&=3.\end{aligned}\)
\end{enumerate}
\end{zadaniReseniBox}
\poznamka{Orientační časové rozvržení: 5 + 12 + 12 + 10 + 6 = 45 min. Přednost má pochopení jádra; rezerva není součástí časového součtu. Domácí úlohy jsou určeny na samostatnou práci, řešení zůstávají učiteli.}
\end{document}
